\documentclass[11pt]{article}
\usepackage[left=1.03in,right=1.03in,top=.88in,bottom=.85in]{geometry}
\usepackage{amsmath,amssymb,amsthm,mathtools}
\usepackage[T1]{fontenc}
\usepackage{lmodern}
\usepackage[expansion=false]{microtype}
\usepackage[hidelinks]{hyperref}
\newtheorem{theorem}{Theorem}
\newtheorem{lemma}[theorem]{Lemma}
\newtheorem{corollary}[theorem]{Corollary}
\newcommand{\eps}{\varepsilon}
\newcommand{\pos}[1]{\bigl(#1\bigr)_{+}}
\newcommand{\Q}{\mathbb{Q}}
\newcommand{\R}{\mathbb{R}}
\newcommand{\one}{\mathbf{1}}
\newcommand{\source}[2]{\href{https://github.com/lsernau/grh-lean/blob/72fead102f663d5c7a407964a13dc0cf9d977207/grh/#1}{#2}}
\setlength{\emergencystretch}{2em}
\title{Sharp sampling of one-sided prime remainders\texorpdfstring{\\}{ }with binary extremal examples}
\author{Research note}
\date{26 November 2026}
\begin{document}
\maketitle
\begin{abstract}
Fourth-power sampling of smoothed prime sums follows from a convex interpolation bound. How far can the sampling grid be thinned if each underlying integer has weight zero or one? We characterize every sampling mesh that preserves an upper decay bound $x^{-\theta+\eps}$ for all $\eps>0$ and all nonnegative weights, for $0<\theta<1$. Necessity already holds for indicator functions of sets of integers of any prescribed density strictly between zero and one. The same classification holds for general error envelopes comparable on dyadic intervals, down to the scale $1/x$. We specialize the interpolation result to the exact twisted prime sum in the Alp\"oge--Sernau reduction to Dirichlet GRH.
\end{abstract}

\section{From zeros to fourth powers}
Write $t_+=\max(t,0)$, let $\Lambda$ be the von Mangoldt function, and let $\varphi$ denote Euler's totient. For a positive integer $m$ put $Q=18480m$. For each prime $d\equiv Q/2-1\pmod Q$ define
\begin{align}
 w_{Q,d}(n)&=\Lambda(n)\one_{n\equiv1\ (Q)}\bigl(1-(n/d)\bigr),\label{eq:w}\\
 R_{Q,d}(x)&=\sum_{n\ge1}w_{Q,d}(n)\pos{1/n-1/x}
       -\frac{\log x}{\varphi(Q)}-C_{Q,d},\qquad x\ge1.\label{eq:R}
\end{align}
Here $(n/d)$ is the Legendre symbol, including its value zero at multiples of $d$. The normalization $C_{Q,d}$ is the limit of the first two terms as $x\to\infty$. Its unconditional existence follows from the fixed-modulus prime number theorem with classical error, or from the normalization theorem in the Lean development~\cite{grh}. In its notation $C_{Q,d}=\Re(\mathrm{logDerivMean}(d,Q,1))$. The factor in parentheses in~\eqref{eq:w} lies in $\{0,1,2\}$.

\begin{theorem}[Fourth-power criterion]\label{t:quartic}
Dirichlet GRH is equivalent to the following assertion. For every $m\ge1$, every prime $d\equiv 9240m-1\pmod{18480m}$, and every $\eps>0$, all sufficiently large positive integers $k$ satisfy
\[
 R_{18480m,d}(k^4)\le (k^4)^{-1/2+\eps}.
\]
The threshold may depend on $m,d,\eps$.
\end{theorem}
Here Dirichlet GRH is the statement that every zero in the open critical strip of every Dirichlet $L$-function has real part $1/2$. The theorem reduces the number of abscissas up to $X$ to $\lfloor X^{1/4}\rfloor$ per twist. Each sampled sum still involves prime powers up to its abscissa. The theorem concerns an infinite family of eventual inequalities.

\section{The interpolation estimate}
Given weights $w(n)\ge0$, a coefficient $c>0$, and $C\in\R$, put
\[
 R(x)=\sum_{n\ge1}w(n)\pos{1/n-1/x}-c\log x-C.
\]
All sums here are finite at finite $x$. The mechanism is that
$P(u)=\sum_n w(n)(1/n-u)_+$ is convex on $(0,\infty)$, while $R(1/u)=P(u)+c\log u-C$.

\begin{lemma}[A bound on a single interval]\label{l:chord}
For $1\le a<b$ and $a\le x\le b$, put $v=1/b$, $z=1/a$, $u=1/x$, and $t=(u-v)/(z-v)$. Then
\begin{align*}
 R(x)&\le (1-t)R(b)+tR(a)+c\{\log u-(1-t)\log v-t\log z\}\\
 &\le \max\{R(a),R(b)\}+\frac{c}{8}\left(\frac{b-a}{a}\right)^2 .
\end{align*}
\end{lemma}
\begin{proof}
The first inequality is the defining chord inequality for $P$. The second derivative of $\log u$ on $[v,z]$ has magnitude at most $1/v^2$. Its excess over its chord is bounded above by $(u-v)(z-u)/(2v^2)$, as follows by subtracting the latter quadratic and applying convexity. The product is at most $(z-v)^2/4$, and $(z-v)/v=(b-a)/a$.
\end{proof}

For an exact finite bound write $r=b/a$ and $L=(r-1)/\log r$. Maximizing $\log y-(y-1)\log r/(r-1)$ over $1\le y\le r$ gives
\[
 \max_{v\le u\le z}\{\log u-(1-t)\log v-t\log z\}=\log L-1+L^{-1}.
\]
This classical scalar quantity is the logarithm of Specht's ratio~\cite{specht}. The quadratic estimate is more convenient on fine meshes.

\begin{proof}[Proof of Theorem~\ref{t:quartic}]
Assume Dirichlet GRH. For fixed modulus $M$, the GRH prime number theorem in a reduced residue class gives
$\psi(x;M,a)=x/\varphi(M)+O_M(\sqrt{x}\log^2(2x))$~\cite{ehp}. Apply it at modulus $Qd$. The congruence defining $d$ gives $(d,Q)=1$ since $Q/2$ is even and $(Q/2-1,Q/2)=1$. Each unit class modulo $d$ occurs once among the unit classes modulo $Qd$ that reduce to $1$ modulo $Q$. The sum of $1-(a/d)$ over these classes is $d-1$. Prime powers divisible by $d$ contribute $O_d(\log(2x))$, and those sharing a prime factor with $Q$ cannot satisfy the progression. Thus, writing $A(x)=\sum_{n\le x}w_{Q,d}(n)$, we obtain
\[
 A(x)=\frac{x}{\varphi(Q)}+O_{Q,d}(\sqrt{x}\log^2(2x)).
\]
Since $\sum_n w(n)(1/n-1/x)_+=\int_1^x A(t)t^{-2}\,dt$, subtracting the limit constant shows that $R_{Q,d}(x)$ is the negative tail integral of $(A(t)-t/\varphi(Q))/t^2$. Its absolute value is $O_{Q,d}(x^{-1/2}\log^2(2x))$. This implies the stated eventual estimate.

Conversely assume the sampled assertion. On $[k^4,(k+1)^4]$, the relative gap is $(1+1/k)^4-1=O(1/k)$. Lemma~\ref{l:chord} incurs $O_Q(k^{-2})=O_Q(x^{-1/2})$. Applying the sampled assertion with $\eps/2$ at the two endpoints gives $R_{Q,d}(x)\le x^{-1/2+\eps}$ eventually for every $\eps>0$, absorbing fixed factors by the spare exponent.

For precision, the analytic input~\cite{grh} applies when $8\mid Q$, $d_0\equiv7\pmod 8$, $(d_0,Q)=1$, and $\beta>1/2$. If every prime $d\equiv d_0\pmod Q$ satisfies $R_{Q,d}(x)\le x^{\beta-1}$ for all sufficiently large real $x$, it yields $L(\rho,\psi)\ne0$ for every $\psi$ modulo $Q$ and $\Re\rho>\beta$, $\rho\ne1$. Here $d_0=9240m-1$ meets these conditions. For a putative zero $1/2<\Re\rho<1$ of a character modulo $m$, choose $1/2<\beta<\Re\rho$ and use $\eps=\beta-1/2$. Induce the character to modulus $Q$. Its $L$-function is the old one times finitely many factors $1-\chi(p)p^{-s}$, so the putative zero persists and contradicts the input. For $0<\Re\rho<1/2$, reduction to a primitive character and the functional equation give a zero on the excluded side for the conjugate primitive character. The extra Euler factors in the critical strip are nonzero.
\end{proof}

\section{Which meshes suffice?}
Let $1\le x_1<x_2<\cdots\to\infty$ and $0<\theta<1$. Say that $R$ satisfies the one-sided rate on a set $S$ if, for every $\eps>0$, there exist $K_\eps,X_\eps>0$ such that
\[
 R(x)\le K_\eps x^{-\theta+\eps}\qquad(x\in S,\ x\ge X_\eps).
\]
The factor $K_\eps$ may equivalently be replaced by $1$ by first using the hypothesis with $\eps/2$.

\begin{theorem}[Classification of sampling meshes]\label{t:mesh}
The following conditions are equivalent.
\begin{enumerate}
\item For each $\delta>0$, eventually
\[
 (x_{j+1}-x_j)/x_j\le x_j^{-\theta/2+\delta}.
\tag{M}
\]
\item For every $w(n)\ge0$, $c>0$ and $C\in\R$, the one-sided rate on $\{x_j\}$ implies that rate on $[1,\infty)$.
\item For any prescribed $c\in(0,1)$, the implication in (2) holds for every indicator $w$ of natural density $c$ for which $R$ is normalized by $R(x)\to0$.
\end{enumerate}
In (3) one may fix a single such $c$ before testing the mesh.
\end{theorem}
\begin{proof}
Under (M) we have $x_{j+1}/x_j\to1$, by taking $\delta<\theta/2$. For $x_j\le x\le x_{j+1}$ the squared relative gap is $O(x^{-\theta+2\delta})$. Fix a target $\eps>0$, use $2\delta<\eps$, and apply Lemma~\ref{l:chord} and the endpoint bounds with exponent $\eps/2$. This gives (2), which implies (3).

To prove necessity, suppose (M) fails. Choose $0<\delta<\theta/2$ so that infinitely many intervals $[a,b]=[x_j,x_{j+1}]$ have
\[
 h=\min\{(b-a)/a,1\}\ge a^{-\theta/2+\delta}.
\]
Use only these intervals with $a\ge3$, and put $H=ah$. Let $f(t)=t^3(1-t)^3$ on $[0,1]$ and zero elsewhere. It is $C^2$, and $f(1/2)>0$. On $[a,a+H]$ define
\[
 B(x)=\eta\,\frac{h^2}{\log a}\,f\!\left(\frac{x-a}{H}\right).
\]
Set $B=0$ on the complement of the chosen subintervals. The zero value and first two zero derivatives at the support endpoints make this a global $C^2$ function, since there are finitely many intervals on each compact set. We have $B(x_j)=0$ for every sample and $B(x)\to0$.

Now prescribe a smooth cumulative count
\[
 F(x)=cx+x^2B'(x),\qquad x\ge0.
\]
On the bump support $x\le2a$ and $H=ah$. There $|2xB'|\le 4\eta h\|f'\|_\infty/\log a$ and $|x^2B''|\le4\eta\|f''\|_\infty/\log a$. A single sufficiently small $\eta>0$ therefore ensures $0\le F'(x)\le1$ for all $x$. Also $F(0)=0$ and $F(x)/x\to c$, using $|xB'|\le2\eta h\|f'\|_\infty/\log a$.

For $n\ge1$ take the binary weights
\[
 w(n)=\lfloor F(n)\rfloor-\lfloor F(n-1)\rfloor.
\]
The derivative bounds imply $w(n)\in\{0,1\}$ and the cumulative $A(t)=\sum_{n\le t}w(n)=\lfloor F(\lfloor t\rfloor)\rfloor$ differs from $F(t)$ by at most $2$. Hence $w$ has density $c$. For $U(x)=\sum_n w(n)(1/n-1/x)_+-c\log x$, integration gives
\[
 U(x)-B(x)=-B(1)+\int_1^x\frac{A(t)-F(t)}{t^2}\,dt .
\]
The right side has a finite limit $C$. The normalized $R=U-C$ satisfies $|R(x)-B(x)|\le2/x$, so $|R(x_j)|\le2/x_j$ at every sample. At $y=a+H/2\asymp a$, however,
\[
 B(y)\ge \eta f(1/2)\,a^{-\theta+2\delta}/\log a .
\]
This dominates both $1/y$ and $y^{-\theta+\delta}$. The full one-sided rate fails for $\eps=\delta$.
\end{proof}

\begin{theorem}[Exact rates, including slowly varying factors]\label{t:gauge}
Let $E(x)>0$ tend to zero. Suppose $1/x=O(E(x))$ and for some $D\ge1$ and all sufficiently large $a$, one has
\[
 D^{-1}E(a)\le E(y)\le DE(a)\qquad(a\le y\le2a).
\]
Write $R=O_+(E)$ for the eventual bound $R(x)\le K E(x)$ with some $K>0$. A mesh preserves $O_+(E)$ from the sampled points to the whole ray for all nonnegative weights exactly when
\[
 \left(\frac{x_{j+1}-x_j}{x_j}\right)^2=O(E(x_j)).
\]
Each mesh failing this condition admits a counterexample with binary weights of any prescribed density $0<c<1$, normalized by $R(x)\to0$.
\end{theorem}
\begin{proof}
Under the stated gap condition $b/a\to1$, so eventually $a\le x\le b\le2a$. Both endpoint error envelopes and the quadratic penalty in Lemma~\ref{l:chord} are $O(E(x))$. Conversely, failure supplies a subsequence with $h^2/E(a)\to\infty$ for $h=\min\{(b-a)/a,1\}$ because $E(a)\to0$. Keep the bump and rounding construction of Theorem~\ref{t:mesh}, replacing $1/\log a$ by $s_a=\sqrt{E(a)}/h\to0$ on this subsequence. Eventually $s_a\le1$, so the same derivative bounds ensure $0\le F'\le1$, $F(x)/x\to c$, $B(x)\to0$, and $|R(x)-B(x)|\le2/x$. Sample errors are $O(E)$, while at each selected midpoint $y$, $B(y)/E(y)\ge (\eta f(1/2)/D)h/\sqrt{E(a)}\to\infty$. Thus $R(y)/E(y)$ is unbounded above as well.
\end{proof}
For example, an upper envelope $x^{-1/2}(\log x)^\alpha$ admits gaps of order $x^{3/4}(\log x)^{\alpha/2}$ for every fixed real $\alpha$. The logarithmic factor is sharp within the class of binary weights as a consequence of Theorem~\ref{t:gauge}.

\begin{corollary}[Superquartic power grids]\label{c:sharp}
For any real $p>4$ and any $0<c<1$, there is a set of positive integers of density $c$ whose normalized remainder satisfies
\[
 |R(a_k)|\le2/a_k,\qquad R((a_k+a_{k+1})/2)\asymp a_k^{-2/p}
\]
for all sufficiently large $k$, where $a_k=\lfloor k^p\rfloor$. The same conclusion holds for the real grid $a_k=k^p$.
\end{corollary}
\begin{proof}
Here $h_k=(a_{k+1}-a_k)/a_k\sim p/k$, since rounding changes the gap by at most one. Use every sufficiently large interval of the grid in the previous construction and replace $1/\log a_k$ by $1$. The same derivative bounds apply with $h_k\le1$, and now $xB'=O(h_k)\to0$ and $B=O(h_k^2)\to0$. The rounding discrepancy remains at most $2/x$ for the normalized remainder. The midpoint bump is exactly $\eta f(1/2)h_k^2\asymp a_k^{-2/p}$; since $2/p<1$, this dominates the rounding error.
\end{proof}
At $\theta=1/2$, mesh (M) permits relative gaps $x^{-1/4+o(1)}$, equivalently absolute gaps $x^{3/4+o(1)}$, as upper bounds. The sharpness statement concerns all positive weights. The proof does not impose the multiplicative or spectral constraints of the actual twisted von Mangoldt weights. Special arithmetic information can improve sampling in related problems, as discussed below.

\section{Sources and scope}
The published repository traces Dirichlet zero-freeness back to a one-sided prime remainder by comparing upper and lower bounds for an index-corrected lattice covolume~\cite{grh}. The interpolation here begins with the positive-part kernel already present in that reduction. The project reports independent checks of its exported theorems and a proof based on Alp\"oge's work with Claude. This note uses its identified analytic implication as an input. The proofs of Lemma~\ref{l:chord}, Theorem~\ref{t:mesh} and Corollary~\ref{c:sharp} are elementary and independent of GRH. The related work in ProveIt is quite close. Setiawan~\cite{setiawan} already gives the fourth-power square-root interpolation scale. The ProveIt report \emph{Sharp interpolation and sparse nonabelian tests}~\cite{prior} already gives the reciprocal chord argument, exact logarithmic defect, and sharp examples supported on power grids whose atom sizes grow like $k^{p-1}$. A second report~\cite{powers} obtains every fixed power grid for full Dedekind prime-ideal weights by using the zero explicit formula. These facts motivate the restriction to weights at most one. Theorems~\ref{t:mesh} and~\ref{t:gauge} establish the matching necessity on arbitrary meshes with this restriction, prescribed natural density, and a vanishing normalized remainder. The fourth-power exponent, the scalar chord bound, and the general connection between one-sided prime sums and zeros have earlier sources~\cite{setiawan,prior,specht,suzuki}. The present comparison does not establish historical priority for the binary converse.

\small
\begin{thebibliography}{9}
\bibitem{grh} L. Sernau and collaborators, \emph{GRH for Dirichlet $L$-functions, in Lean 4}, revision \texttt{72fead102f66}, 25 November 2026.
See the \source{README.md}{project verification report}, \source{src/Main.lean}{exported theorems}, \source{src/New/Reduction/Zeros/PrimeError.lean}{weights and normalization}, \source{src/New/Reduction/Zeros/Criterion.lean}{one-sided analytic criterion}, \source{src/New/Reduction/MainInequality.lean}{covolume comparison}, and \source{src/New/Reduction/OneSidedBound.lean}{specialization and induction of characters}.
\bibitem{ehp} A.-M. Ernvall-Hyt\"onen and N. Paloj\"arvi, \emph{Explicit bound for the number of primes in arithmetic progressions assuming the Generalized Riemann Hypothesis}. \href{https://arxiv.org/abs/2003.01925}{arXiv:2003.01925}.
\bibitem{suzuki} M. Suzuki, \emph{On variants of Chebyshev's conjecture}. \href{https://arxiv.org/abs/2411.07436}{arXiv:2411.07436}. \href{https://doi.org/10.1007/s11139-025-01289-y}{Published correction}.
\bibitem{specht} M. Tominaga, \emph{Specht's ratio and logarithmic mean in the Young inequality}, Math. Inequal. Appl. \textbf{7} (2004), 113--125. \href{https://files.ele-math.com/articles/mia-07-13.pdf}{Paper}.
\bibitem{setiawan} S. Setiawan, \emph{The Primacohedron trace formulation of the Riemann hypothesis}, version 2 (2026). \href{https://www.preprints.org/manuscript/202611.0738}{Preprint}.
\bibitem{prior} ProveIt, \emph{Sharp interpolation and sparse nonabelian tests}, in \emph{Beyond Dirichlet GRH}, November 2026. \href{https://github.com/VladimirReshetnikov/ProveIt/blob/1f20b04e3f8d78f16d1a873547d90554e46e40bf/Analysis/GeneralizedRiemannHypothesis/Research/grh_method_transfers_2026-11-25/report/sparse_sampling.tex}{Archived section}.
\bibitem{powers} ProveIt, \emph{What sparse prime sums determine about zeta zeros}, November 2026. \href{https://github.com/VladimirReshetnikov/ProveIt/blob/1f20b04e3f8d78f16d1a873547d90554e46e40bf/Analysis/GeneralizedRiemannHypothesis/Research/grh-torsion-and-prime-sampling/sections/powers.tex}{Section on arbitrary power grids}.
\end{thebibliography}
\end{document}
